% Section 6 - Financial Projection (Financial Forecasting)
\section{Financial Projection}\label{sec:financial}
\emph{Three-year financial plan and financial feasibility assessment.}

%====================================================================
\subsection{Overview and Basis of Financial Projections}
\begin{itemize}[itemsep=0pt]
  \item Uvie is an interdisciplinary product combining chemical engineering and
    software technology. The patch uses UV-sensitive photochromic pigments
    integrated into a skin-compatible substrate. Users apply sunscreen to both
    their skin and the patch. As the sunscreen protection decreases, UV exposure
    causes the patch to change color. The Uvie mobile application uses the phone
    camera to analyze the color intensity and sends a notification reminding
    users to reapply sunscreen.
  \item For commercial planning, this report assumes that Uvie generates revenue
    primarily by selling patch packages. The basic scanning and reminder
    features of the mobile application are assumed to be free of charge. This
    approach is intended to reduce adoption barriers, encourage repeated use and
    support repeat purchases. Revenue from premium subscriptions, advertising,
    business partnerships and exports is excluded to avoid overstating the
    project's potential.
\end{itemize}

\subsubsection{Key Financial Assumptions}
\begin{table}[ht]
\centering\small
\caption{Key financial assumptions (base case)}\label{tab:assumptions}
\begin{tabular}{@{}L{4.0cm}L{3.2cm}L{7.0cm}@{}}
  \toprule
  \textbf{Assumption} & \textbf{Base-case value} & \textbf{Explanation}\\
  \midrule
  Average selling price & 89,000 VND/package & Assumed retail price for one Uvie patch package\\
  Direct cost per package & 34,000 VND/package & Materials, pigments, processing, packaging and quality inspection\\
  Year 1 sales volume & 18,000 packages & Launch stage, mainly through online sales\\
  Year 2 sales volume & 50,000 packages & Expansion of brand awareness and distribution channels\\
  Year 3 sales volume & 100,000 packages & Wider market coverage and higher repeat-purchase volume\\
  Corporate income tax & 20\% & Assumed standard rate; confirm applicable tax treatment before use\\
  NPV discount rate & 18\% per year & Assumed required return for an early-stage startup\\
  \bottomrule
\end{tabular}
\end{table}

\subsubsection{Modeling Principles and Limitations}
\begin{itemize}[itemsep=0pt]
  \item The model uses a bottom-up approach: sales volume multiplied by average
    selling price determines revenue; sales volume multiplied by unit cost
    determines cost of goods sold.
  \item Operating expenses are estimated based on team size, marketing
    activities, product research and development, and maintenance of the mobile
    application.
  \item No revenue from advertising, premium app subscriptions, B2B contracts or
    exports is included.
  \item Inflation, foreign exchange movements, complex financing costs and
    changes in tax regulations are not modeled.
\end{itemize}

%====================================================================
\subsection{Revenue and Cost Forecast}

\subsubsection{Revenue Forecast}
Revenue is calculated by multiplying the number of packages sold by the average
selling price. The projected sales growth assumes that Uvie completes its MVP,
validates the accuracy of color recognition, builds user trust and expands its
distribution channels.

\[
  \text{Net Revenue} = \text{Number of Packages Sold} \times \text{Average Selling Price}
\]
For Year 1: $18{,}000 \times 89{,}000\ \text{VND} = 1{,}602{,}000{,}000\ \text{VND}$,
equivalent to VND 1,602 million.

\begin{table}[ht]
\centering\small
\caption{Revenue forecast}\label{tab:revenue}
\begin{tabular}{@{}L{6.2cm}rrr@{}}
  \toprule
  \textbf{Metric} & \textbf{Year 1} & \textbf{Year 2} & \textbf{Year 3}\\
  \midrule
  Packages sold & 18,000 & 50,000 & 100,000\\
  Average selling price (VND/package) & 89,000 & 89,000 & 89,000\\
  Revenue (VND million) & 1,602 & 4,450 & 8,900\\
  Revenue growth & -- & 177.8\% & 100.0\%\\
  \bottomrule
\end{tabular}
\end{table}

\subsubsection{Cost of Goods Sold (COGS)}
COGS includes UV-sensitive pigments, hydrocolloid/silicone substrates, coating
materials, packaging, processing and quality inspection. The assumed average cost
is VND 34,000 per package at small-to-medium production scale.

\begin{table}[ht]
\centering\small
\caption{Cost of goods sold and gross profit}\label{tab:cogs}
\begin{tabular}{@{}L{6.2cm}rrr@{}}
  \toprule
  \textbf{Metric (VND million)} & \textbf{Year 1} & \textbf{Year 2} & \textbf{Year 3}\\
  \midrule
  Sales volume (packages) & 18,000 & 50,000 & 100,000\\
  Unit cost (VND/package) & 34,000 & 34,000 & 34,000\\
  Total COGS & 612 & 1,700 & 3,400\\
  Gross profit & 990 & 2,750 & 5,500\\
  Gross profit margin & 61.8\% & 61.8\% & 61.8\%\\
  \bottomrule
\end{tabular}
\end{table}

\subsubsection{Operating Expense Forecast}
\begin{table}[ht]
\centering\small
\caption{Operating expense forecast}\label{tab:opex}
\begin{tabular}{@{}L{7.4cm}rrr@{}}
  \toprule
  \textbf{Expense category (VND million)} & \textbf{Year 1} & \textbf{Year 2} & \textbf{Year 3}\\
  \midrule
  Personnel and team operations & 540 & 780 & 1,080\\
  Marketing and sales & 360 & 600 & 1,000\\
  Research and development (R\&D) & 180 & 240 & 300\\
  Administration, legal and testing & 150 & 210 & 300\\
  Application maintenance and digital infrastructure & 90 & 120 & 180\\
  \midrule
  \textbf{Total operating expenses before depreciation} & \textbf{1,320} & \textbf{1,950} & \textbf{2,860}\\
  Depreciation & 90 & 110 & 140\\
  \bottomrule
\end{tabular}
\end{table}

Testing and legal costs should be updated based on the product's regulatory
classification, requirements for skin-contact materials, applicable product
notification/labeling rules and advertising regulations. The estimates above are
not supplier quotations or legal conclusions.

%====================================================================
\subsection{Projected Income Statement}
The projected income statement illustrates Uvie's ability to generate profit over
the first three years. Unit: VND million.

\begin{table}[ht]
\centering\small
\caption{Projected income statement (VND million)}\label{tab:income}
\begin{tabular}{@{}L{8.0cm}rrr@{}}
  \toprule
  \textbf{Income statement item} & \textbf{Year 1} & \textbf{Year 2} & \textbf{Year 3}\\
  \midrule
  Net revenue & 1,602 & 4,450 & 8,900\\
  (--) Cost of goods sold & (612) & (1,700) & (3,400)\\
  \textbf{Gross profit} & \textbf{990} & \textbf{2,750} & \textbf{5,500}\\
  (--) Operating expenses before depreciation & (1,320) & (1,950) & (2,860)\\
  \textbf{EBITDA} & \textbf{(330)} & \textbf{800} & \textbf{2,640}\\
  (--) Depreciation & (90) & (110) & (140)\\
  \textbf{Earnings before interest and tax (EBIT)} & \textbf{(420)} & \textbf{690} & \textbf{2,500}\\
  (--) Interest expense & (30) & (25) & (15)\\
  \textbf{Earnings before tax (EBT)} & \textbf{(450)} & \textbf{665} & \textbf{2,485}\\
  (--) Estimated corporate income tax & 0 & (133) & (497)\\
  \midrule
  \textbf{Net profit after tax} & \textbf{(450)} & \textbf{532} & \textbf{1,988}\\
  \bottomrule
\end{tabular}
\end{table}

\subsubsection{Income Statement Analysis}
\begin{itemize}[itemsep=0pt]
  \item \textbf{Year 1:} Uvie is expected to incur a loss because personnel,
    product development and marketing costs arise before sales reach a
    sufficient scale. This is common during the commercialization stage.
  \item \textbf{Year 2:} EBITDA is projected to turn positive at VND 800 million,
    with estimated net profit after tax of VND 532 million. This depends heavily
    on achieving 50,000 package sales and maintaining the assumed unit cost.
  \item \textbf{Year 3:} Net profit after tax is projected to increase to
    VND 1,988 million as sales volume grows and economies of scale are achieved.
    This outcome requires stable product quality and controlled customer
    acquisition costs.
\end{itemize}

\subsubsection{Break-even Analysis}
\begin{itemize}[itemsep=0pt]
  \item Contribution per package = Selling price $-$ Direct variable cost
    $= 89{,}000 - 34{,}000 = 55{,}000$ VND.
  \item If Year 1 operating expenses and selling costs before depreciation of
    VND 1,320 million are treated as fixed costs for a simplified calculation,
    the operating break-even volume is approximately
    \[
      \frac{1{,}320{,}000{,}000\ \text{VND}}{55{,}000\ \text{VND}} = 24{,}000
      \ \text{packages per year}.
    \]
    This is higher than the projected Year 1 sales volume of 18,000 packages;
    therefore, Uvie is not expected to reach operating break-even in its first
    year.
  \item This simplified calculation treats all operating expenses as fixed. A
    detailed model should separate fixed and variable costs and include sales
    discounts, returns, marketplace fees and distributor commissions.
\end{itemize}

%====================================================================
\subsection{Projected Balance Sheet and Funding Requirements}
The following simplified projected balance sheet illustrates the structure of
assets, liabilities and shareholders' equity. Cash, working capital and
liabilities should ultimately be linked to a detailed cash flow forecast,
collection and payment schedules, inventory assumptions and debt repayment plan.

\subsubsection{Initial Funding Requirements}
\begin{table}[ht]
\centering\small
\caption{Initial funding requirements}\label{tab:funding}
\begin{tabular}{@{}L{4.4cm}L{2.6cm}L{7.2cm}@{}}
  \toprule
  \textbf{Use of funds} & \textbf{Amount (VND million)} & \textbf{Description}\\
  \midrule
  Equipment, tools and fixed assets & 450 & Testing equipment, inspection tools and assets for prototype production\\
  Commercialization preparation & 150 & Packaging design, initial testing, legal preparation and sales-channel setup\\
  Initial working capital & 300 & Raw material inventory, finished goods, cash and initial operating expenses\\
  \midrule
  \textbf{Total initial funding requirement} & \textbf{900} & Assumed founders' equity/investor funding\\
  \bottomrule
\end{tabular}
\end{table}

\subsubsection{Illustrative Projected Balance Sheet}
\begin{table}[ht]
\centering\small
\caption{Illustrative projected balance sheet (VND million)}\label{tab:balance}
\begin{tabular}{@{}L{6.4cm}rrrr@{}}
  \toprule
  \textbf{Balance sheet item} & \textbf{Opening} & \textbf{End Y1} & \textbf{End Y2} & \textbf{End Y3}\\
  \midrule
  \textbf{Assets} & & & &\\
  \quad Cash and cash equivalents$^{*}$ & 300 & 120 & 420 & 1,650\\
  \quad Inventory and net receivables$^{*}$ & 300 & 430 & 650 & 1,000\\
  \quad Net property, plant and equipment & 450 & 460 & 450 & 460\\
  \textbf{Total assets} & \textbf{1,050} & \textbf{1,010} & \textbf{1,520} & \textbf{3,110}\\
  \midrule
  \textbf{Liabilities} & & & &\\
  \quad Short/long-term borrowings and payables$^{*}$ & 150 & 560 & 538 & 140\\
  \textbf{Shareholders' equity} & & & &\\
  \quad Initial contributed capital & 900 & 900 & 900 & 900\\
  \quad Accumulated retained earnings & 0 & (450) & 82 & 2,070\\
  \textbf{Total liabilities and equity} & \textbf{1,050} & \textbf{1,010} & \textbf{1,520} & \textbf{3,110}\\
  \bottomrule
\end{tabular}
\end{table}

{\small $^{*}$ Cash, inventory/net receivables and liabilities are illustrative
balances designed to demonstrate the balance sheet structure; they have not been
fully derived from an integrated cash flow statement. Opening resources comprise
VND 900 million of contributed capital and an assumed VND 150 million of
payables/short-term borrowing. Retained earnings are accumulated from the
projected net profit after tax.\par}

Before submission as a formal financial statement, this table should be replaced
with a fully linked balance sheet based on the cash flow statement, depreciation
schedule, inventory, receivables, payables and debt repayment schedule. The
illustrative balance sheet should not be presented as verified accounting data.

\subsubsection{Potential Funding Sources}
\begin{itemize}[itemsep=0pt]
  \item Founders' capital and angel investors: finance equipment, formulation
    refinement and product testing.
  \item Innovation competitions and incubator programs: support testing, user
    trials and MVP development.
  \item Contract manufacturing and small-batch production: reduce machinery
    investment, inventory exposure and capacity risk.
  \item Scale production only after validating color stability, material safety,
    defect rates and customer repeat-purchase behavior.
\end{itemize}

%====================================================================
\subsection{Financial Ratios and Customer Economics}

\subsubsection{CAC/COCA -- Customer Acquisition Cost}
In this report, COCA is interpreted as Cost of Customer Acquisition, equivalent to
CAC: the cost of acquiring one new customer. It is calculated by dividing
marketing and sales expenditure by the number of new customers acquired during
the period.

\begin{table}[ht]
\centering\small
\caption{Customer acquisition cost (CAC/COCA)}\label{tab:cac}
\begin{tabular}{@{}L{6.2cm}rrr@{}}
  \toprule
  \textbf{Metric} & \textbf{Year 1} & \textbf{Year 2} & \textbf{Year 3}\\
  \midrule
  Marketing and sales cost (VND million) & 360 & 600 & 1,000\\
  New customers (people) & 12,000 & 30,000 & 60,000\\
  CAC/COCA (VND/customer) & 30,000 & 20,000 & 16,667\\
  \bottomrule
\end{tabular}
\end{table}

CAC is assumed to decline over time due to stronger brand awareness,
word-of-mouth referrals and improved marketing efficiency. This assumption must be
validated using campaign-level data. Actual CAC may be higher if discounts,
commissions and sales personnel costs are included.

\subsubsection{LTV -- Customer Lifetime Value}
LTV is estimated as the gross profit generated by one customer over the period in
which they continue purchasing the product. Assuming that one customer purchases
an average of three packages per year, remains active for two years, and generates
gross profit of VND 55,000 per package:
\[
  \text{LTV} = 55{,}000\ \text{VND} \times 3\ \text{packages/year} \times 2\ \text{years}
  = 330{,}000\ \text{VND per customer}.
\]
This simplified formula does not discount future cash flows and does not deduct
customer service, return, or application operating costs. LTV should be updated
using actual cohort retention and repeat-purchase data.

\subsubsection{LTV/CAC Ratio}
\begin{table}[ht]
\centering\small
\caption{LTV/CAC ratio}\label{tab:ltvcac}
\begin{tabular}{@{}L{6.2cm}rrr@{}}
  \toprule
  \textbf{Metric} & \textbf{Year 1} & \textbf{Year 2} & \textbf{Year 3}\\
  \midrule
  LTV (VND/customer) & 330,000 & 330,000 & 330,000\\
  CAC (VND/customer) & 30,000 & 20,000 & 16,667\\
  LTV/CAC & $11.0\times$ & $16.5\times$ & $19.8\times$\\
  \bottomrule
\end{tabular}
\end{table}

An LTV/CAC ratio above 1 indicates that estimated lifetime gross profit exceeds
the cost of acquiring a customer. However, ratios of $11.0\times$--$19.8\times$ are
highly attractive and may reflect optimistic assumptions. The team should validate
purchase frequency, retention, fully loaded CAC and actual gross margin. This
ratio alone is not sufficient to conclude that the project is financially viable.

\subsubsection{Additional Performance Indicators}
\begin{itemize}[itemsep=0pt]
  \item Gross profit margin $= (\text{Revenue} - \text{COGS}) / \text{Revenue}
    = 61.8\%$ under the base-case assumptions.
  \item EBITDA margin is approximately $-20.6\%$, $18.0\%$ and $29.7\%$ in
    Years 1, 2 and 3, respectively.
  \item Repeat-purchase rate, return rate, defective patch rate and
    warranty/support costs should be monitored as key operating KPIs.
\end{itemize}

%====================================================================
\subsection{NPV, IRR and Financial Feasibility Assessment}
Net Present Value (NPV) measures the present value of a project's expected cash
flows after discounting them at the required rate of return. Internal Rate of
Return (IRR) is the discount rate at which NPV equals zero.

\subsubsection{Illustrative Project Cash Flows}
For illustration, the model assumes an initial investment of VND 900 million and
the following estimated free cash flow to the firm (FCFF). FCFF is approximated as
after-tax EBIT plus depreciation, less capital expenditure and increases in net
working capital.

\begin{table}[ht]
\centering\small
\caption{Illustrative free cash flow to the firm (VND million)}\label{tab:fcff}
\begin{tabular}{@{}L{6.0cm}rrrr@{}}
  \toprule
  \textbf{Cash flow item} & \textbf{Year 0} & \textbf{Year 1} & \textbf{Year 2} & \textbf{Year 3}\\
  \midrule
  Initial investment & (900) & -- & -- & --\\
  EBIT & -- & (420) & 690 & 2,500\\
  Tax on EBIT (estimated) & -- & 0 & (138) & (500)\\
  Add: depreciation & -- & 90 & 110 & 140\\
  Capital expenditure (CAPEX) & -- & (100) & (100) & (150)\\
  Increase in net working capital & -- & (100) & (150) & (200)\\
  \midrule
  \textbf{Free cash flow (FCFF)} & \textbf{(900)} & \textbf{(530)} & \textbf{412} & \textbf{1,790}\\
  \bottomrule
\end{tabular}
\end{table}

\subsubsection{NPV and IRR Results}
Using an assumed annual discount rate of 18\%:
\[
  \text{NPV} = \sum_{t=0}^{3} \frac{\text{FCFF}_t}{(1 + 18\%)^{t}}
  \approx +36.2\ \text{VND million},
  \qquad
  \text{IRR} \approx 19.3\%.
\]
Under the base-case scenario, NPV is positive and IRR is slightly higher than the
18\% discount rate, indicating that the project marginally exceeds the required
return. However, the margin of safety is limited: NPV is only approximately
VND 36.2 million. A relatively small decline in sales volume or increase in unit
cost or marketing expenditure could make NPV negative.

\subsubsection{Sensitivity and Risk Analysis}
\begin{table}[ht]
\centering\small
\caption{Key risk factors and mitigation}\label{tab:risk}
\begin{tabular}{@{}L{4.0cm}L{5.0cm}L{5.0cm}@{}}
  \toprule
  \textbf{Risk factor} & \textbf{Potential financial impact} & \textbf{Mitigation}\\
  \midrule
  Sales volume below forecast & Lower revenue and longer time to break even & Run small-batch market tests and track conversion rates\\
  Increase in unit production cost & Lower gross profit margin & Maintain 2--3 suppliers; optimize materials and packaging\\
  Higher customer acquisition cost & Lower LTV/CAC and weaker cash flow & Measure CAC by channel; prioritize referrals and educational content\\
  Unstable sensor accuracy & More returns, lower repeat purchases and reputational damage & Test under different UV, sweat and lighting conditions\\
  Additional testing or regulatory requirements & Higher costs and delayed launch & Set aside a testing budget and obtain professional advice\\
  \bottomrule
\end{tabular}
\end{table}

%====================================================================
\subsection{Conclusion and Recommendations}
\begin{itemize}[itemsep=0pt]
  \item Based on the illustrative assumptions, Uvie has the potential to achieve a
    gross profit margin of approximately 61.8\%, turn EBITDA positive in Year 2
    and generate estimated net profit after tax of VND 1,988 million in Year 3.
    Estimated NPV is positive at VND 36.2 million, while IRR is approximately
    19.3\%, slightly above the assumed 18\% discount rate. These results suggest
    that the project may be financially feasible under the base-case scenario,
    although its margin of safety is relatively narrow.
  \item Before finalizing the report, the team should obtain quotations for patch
    materials and manufacturing, confirm the number of patches per package, test
    customers' willingness to pay, develop channel-specific marketing budgets and
    conduct pilot sales. In particular, the safety of skin-contact materials and
    the reliability of the color-recognition algorithm must be validated. These
    factors directly affect unit cost, testing expenses, repeat-purchase rates and
    cash flow.
  \item \textbf{Final conclusion:} The figures presented are scenario-based
    projections, not guaranteed returns. The team should update the model using
    actual data and prepare conservative, base-case and optimistic scenarios
    before making investment decisions.
\end{itemize}
